\section{Nội dung lý thuyết}

\subsection{Kí hiệu và quy ước toán học}

Chúng ta bắt đầu với một không gian mẫu $\Omega$ được trang bị một $\sigma$-đại
số $\mathcal{F}$, và tập hợp $\mathcal{M}_1$ gồm tất cả các độ đo xác suất trên
$(\Omega, \mathcal{F})$. Các phần tử của $\mathcal{M}_1$ còn được gọi là các
phân phối. Theo định nghĩa hình thức, một biến ngẫu nhiên là một hàm đo được từ
$(\Omega, \mathcal{F})$ đến $\mathbb{R}$ hoặc $[-\infty, \infty]$. Khi cho phép
một biến ngẫu nhiên nhận các giá trị vô hạn, chúng ta sẽ ghi chú rõ ràng; nếu
không, nó sẽ nhận giá trị trong $\mathbb{R}$. Chúng ta giả định rằng một phân
phối nào đó $P^\dagger \in \mathcal{M}_1$ chi phối dữ liệu $X$ của chúng ta (là
một biến ngẫu nhiên). Lưu ý rằng $X$ có thể là một véc-tơ $(X_1, \dots, X_n)$.
Các biến $X_1, \dots, X_n$ có thể độc lập và cùng phân phối (iid) dưới
$P^\dagger$, nhưng điều này không nhất thiết phải xảy ra. Chúng ta sử dụng
$X^t$ như một cách viết tắt cho $t$ điểm dữ liệu đầu tiên $(X_1, \dots, X_t)$.

\begin{defn}[Giả thuyết]
  Một giả thuyết (hypothesis) là một tập hợp các độ đo xác suất trong
  $\mathcal{M}_1$. Một giả thuyết là đơn (single) nếu nó chỉ là một tập đơn tử,
  chẳng hạn như $\{P\}$ và $\{Q\}$. Ngược lại, nó được gọi là giả thuyết hợp
  (composite).
\end{defn}

Chúng ta sẽ luôn sử dụng $\mathcal{P}$ cho giả thuyết không và $\mathcal{Q}$
cho giả thuyết thay thế. Chúng ta cũng sẽ viết $H_0$ là giả thuyết không rằng
$P^\dagger \in \mathcal{P}$, và $H_1$ là giả thuyết thay thế rằng $P^\dagger
\in \mathcal{Q}$; nghĩa là, chúng ta hiểu một giả thuyết vừa như một đối tượng
toán học $\mathcal{P}$, vừa như một phát biểu thống kê $H_0$. Đối với một giả
thuyết đơn $\mathcal{P} = \{P\}$, chúng ta sử dụng các cụm từ ``kiểm định
$H_0$'', ``kiểm định $\mathcal{P}$'' và ``kiểm định $P$'' thay thế cho nhau.
Các tập hợp $\mathcal{P}$ và $\mathcal{Q}$ luôn được giả định là rời nhau. Gọi
$\text{Conv}(\mathcal{P})$ là bao lồi của $\mathcal{P}$.

Như đã đề cập ở trên, chúng ta luôn sử dụng $\mathcal{P}, \mathcal{Q}$ cho các
tập các phân phối, và $P, Q$ cho một phân phối đơn lẻ. Trong bối cảnh của một
giả thuyết $\mathcal{P}$, các phát biểu về tính độc lập, phân phối và đẳng thức
hầu chắc chắn (\textit{almost-sure}) được hiểu là đúng đối với mọi phần tử
trong tập đó. Chúng ta nói rằng $\mathcal{Q}$ liên tục tuyệt đối đối với
$\mathcal{P}$, ký hiệu $\mathcal{Q} \ll \mathcal{P}$, nếu $P(A) = 0$ kéo theo
$Q(A) = 0$ với mọi tập đo được $A$. Trong trường hợp này, chúng ta sử dụng
$dQ/dP$ để chỉ \gls{likelihood-ratio} (hoặc tổng quát hơn
là đạo hàm Radon–Nikodym) giữa $Q$ và $P$. Nếu ký hiệu $q$ và $p$ lần lượt là
các hàm mật độ của $Q$ và $P$ đối với một độ đo tham chiếu $L$ (thường là độ đo
Lebesgue, do đó mới có ký hiệu này), chúng ta có thể viết $(dQ/dP)(x) =
q(x)/p(x)$. Kỳ vọng của một biến ngẫu nhiên $X$ dưới $P$ được ký hiệu là
$\expect_P[X] = \int X\,dP$.

Một số ký hiệu sẽ được sử dụng xuyên suốt. Tập $\mathbb{N} = \{1, 2, \dots \}$
là tập tất cả các số nguyên dương, và $\mathbb{R} = (-\infty, \infty)$ là tập
tất cả các số thực. Chúng ta nói rằng một biến ngẫu nhiên $X$ lớn hơn một cách
ngẫu nhiên (dưới $P$, đôi khi được lược bỏ nếu ngữ cảnh đã rõ) so với một phân
phối trên $\mathbb{R}$, được biểu diễn bởi hàm phân phối tích lũy (cdf) $F$,
nếu $P(X \leq x) \leq F(x)$ với mọi $x \in \mathbb{R}$. Chúng ta viết $X
\stackrel{d}{\sim} F$ nếu phân phối của $X$ là $F$ dưới một độ đo xác suất $P$
(ngữ cảnh phải rõ ràng); ký hiệu $\stackrel{iid}{\sim}$ cũng tương tự. Với một
số nguyên dương $K$, chúng ta ký hiệu $[K] = \{1, \dots, K\}$ và $\Delta_K$ là
đơn hình chuẩn trong $\mathbb{R}^K$, tức là:
\begin{equation}
    \Delta_K = \biggl\{ (v_1, \dots, v_K) \in [0, 1]^K : \sum_{k=1}^K v_k = 1
    \biggr\}.
\end{equation} Số Euler $\exp(1)$ được ký hiệu là $e$ (xin lưu ý rằng điều này
khác với biến $e$, nhưng tất nhiên chúng ta sẽ tránh sử dụng cả hai trong cùng
một phương trình bất cứ khi nào có thể).

Tất cả các thuật ngữ như ``tăng'' và ``giảm'' đều được hiểu theo nghĩa không
ngặt.  Chúng ta viết $x \wedge y$ và $x \vee y$ lần lượt cho giá trị nhỏ nhất
và giá trị lớn nhất của hai số thực $x, y$.

\subsection{e-variable, p-variable và test}

\begin{defn}
  Cho $\mathcal{P} \subseteq \mathcal{M}_1$ là một tập hợp các độ đo xác suất
  đại diện cho giả thuyết không. \begin{enumerate}[label=(\roman*)]
    \item Một e-variable $E$ cho $\mathcal{P}$ là một biến ngẫu nhiên nhận
      giá trị trong $[0, \infty]$ thỏa mãn $\expect_P[E] \leq 1$ với mọi $P \in
      \mathcal{P}$. Một e-variable $E$ được gọi là chính xác (\emph{exact})
      nếu $\expect_P[E] = 1$ với mọi $P \in \mathcal{P}$.

    \item Một p-variable $P$ cho $\mathcal{P}$ là một biến ngẫu nhiên nhận giá
      trị trong $[0, 1]$ thỏa mãn $\prob(P \leq \alpha) \leq \alpha$ với mọi
      $\alpha \in (0, 1)$ và mọi $P \in \mathcal{P}$. Một p-variable $P$ được
      gọi là chính xác nếu $\prob(P \leq \alpha) = \alpha$ với mọi $\alpha \in
      (0, 1)$ và mọi $P \in \mathcal{P}$.

    \item Một \gls{test} $\phi$ là một biến ngẫu nhiên nhận giá trị trong
      $[0, 1]$. Một kiểm định được gọi là nhị phân (\emph{binary}) nếu phạm vi
      giá trị của nó là $\{0, 1\}$. Tỉ lệ sai lầm loại I của một kiểm định
      $\phi$ đối với $P$ là $\expect_P[\phi]$. Một kiểm định $\phi$
      có mức ý nghĩa $\alpha \in [0, 1]$ đối với $\mathcal{P}$ nếu tỉ lệ sai
      lầm loại I của nó không vượt quá $\alpha$ với mỗi $P \in \mathcal{P}$.
  \end{enumerate}
\end{defn}

Giá trị mà e-variable hoặc p-variable nhận được sau khi quan sát dữ liệu được
gọi là e-value hoặc p-value. Do đó, có thể thấy rằng e/p-values thực chất không
phải là các đối tượng được định nghĩa chặt chẽ về mặt toán học, mặc dù đây là
những thuật ngữ phổ biến trong các tài liệu thống kê. Trong cuốn sách này,
chúng tôi luôn sử dụng thuật ngữ ``e/p-variable'' khi đưa ra các phát biểu hoặc
kết quả toán học về các biến ngẫu nhiên tương ứng, và đề cập đến ``e/p-value''
khi ý nghĩa thống kê của chúng quan trọng hơn. Trong trường hợp sau, e/p-value
có thể chỉ biến ngẫu nhiên hoặc giá trị thực hiện của nó, và điều này sẽ được
làm rõ dựa trên ngữ cảnh.

\subsection{Test supermartingale, sequential test, and e-process}

Có bốn công cụ SAVI chính: e-process, p-process, $(1 - \alpha)$-\gls{cs} và
\gls{st} mức ý nghĩa $\alpha$; dễ dàng nhận thấy rằng hai công
cụ sau gắn liền với một mức ý nghĩa $\alpha \in [0, 1]$ được xác định trước,
trong khi hai công cụ đầu thì không. Đây lần lượt là các phiên bản tuần tự của
e-variable, p-variable, $(1 - \alpha)$-\gls{cs} và \gls{test} mức ý nghĩa
$\alpha$. Chúng ta sẽ tập trung chủ yếu vào e-process vì chúng là đối tượng
trung tâm; các đối tượng còn lại có thể dễ dàng suy ra từ bất đẳng thức Ville
áp dụng cho một hoặc nhiều e-process.

Chúng ta bắt đầu với nền tảng kỹ thuật cần thiết cho bối cảnh tuần tự. Chúng ta
giả định rằng người đọc đã quen thuộc với các thuật ngữ trong tiến trình ngẫu
nhiên, được tóm tắt dưới đây.

Chúng ta làm việc trong một không gian xác suất \gls{filtered} cố định,
nghĩa là $\mathcal{F}$ đề cập đến một \gls{filtration} $(\mathcal{F}_t)_{t
\in T}$ (một dãy các $\sigma$-đại số lồng nhau), trong đó $T$ là một tập hợp
chỉ số hữu hạn hoặc vô hạn bắt đầu từ 0. Chúng ta chỉ xử lý các tiến trình thời
gian rời rạc, nên $T = \mathbb{N}_0$ hay $\{0, 1, \dots, n\}$ với $n \in
\mathbb{N}$. Gọi $\mathcal{F}_\infty = \bigcup_{t \in T} \mathcal{F}_t$. Ký
hiệu $\mathcal{T}_+$là tập hợp $T \setminus \{0\}$.

Một dãy các biến ngẫu nhiên $Y = (Y_t)_{t \in T}$ được gọi là một tiến trình
nếu nó thích nghi với $\mathcal{F}$—tức là, nếu $Y_t$ đo được đối với
$\mathcal{F}_t$ với mọi $t$. Một tiến trình $Y$ được gọi là có thể dự báo nếu
$Y_t$ đo được đối với $\mathcal{F}_{t-1}$ với mỗi $t \geq 1$. Các bất đẳng thức
so sánh các tiến trình được hiểu là giữ nguyên tại mỗi $t \in T$.

Thông thường $\mathcal{F}$ được chọn là \gls{filtration} tự nhiên của dữ liệu,
đó là $\mathcal{F}_t := \sigma(X^t)$, trong đó $X^t = (X_1, \dots, X_t)$, với
$\mathcal{F}_0$ là tầm thường ($\mathcal{F}_0 = \{\emptyset, \Omega\}$). Trong
trường hợp này, việc $Y_t$ đo được đối với $\mathcal{F}_t$ có nghĩa là $Y_t$ là
một hàm đo được của $X^t$. Tuy nhiên, đôi khi $\mathcal{F}$ là một
\gls{filtration} thô hơn (coarser) (chúng ta loại bỏ thông tin) hoặc phong phú
hơn (richer) (ví dụ cho phép $\mathcal{F}_0 = \sigma(U_0)$ đối với một biến
ngẫu nhiên phân phối đều $U_0$ độc lập với tất cả các yếu tố ngẫu nhiên khác).

Một thời điểm (quy tắc) dừng $\tau$ là một biến ngẫu nhiên nhận giá trị nguyên
không âm sao cho $\{\tau \leq t\} \in \mathcal{F}_t$ với mỗi $t \geq 0$. Nói
cách khác: tại mỗi thời điểm, chúng ta biết liệu quy tắc dừng đang ra lệnh cho
chúng ta dừng lại hay tiếp tục. Ký hiệu $\mathcal{T}$ là tập hợp tất cả các
thời điểm dừng, phụ thuộc ngầm vào $\mathcal{F}$, bao gồm cả những thời điểm có
thể không bao giờ dừng. Đối với một thời điểm dừng $\tau$, $\sigma$-đại
số $\mathcal{F}_\tau$ được định nghĩa bởi $\{A \in \mathcal{F}_\infty : A \cap
\{\tau \leq t\} \in \mathcal{F}_t \text{ với mọi } t \in T\}$.

Ta biết rằng, tập hợp các phân phối $\mathcal{P}$ đại diện cho giả thuyết
không. Một \gls{st} mức ý nghĩa $\alpha$ cho $\mathcal{P}$ là một tiến trình
nhị phân $\phi$ sao cho:
\begin{equation}
    \sup_{P \in \mathcal{P}} \prob_P(\exists t \geq 1 : \phi_t = 1) \leq
    \alpha.
\end{equation}

Giải thích $\phi_t = 1$ là bác bỏ $\mathcal{P}$, điều kiện trên có nghĩa là xác
suất mà chúng ta sẽ bác bỏ sai giả thuyết không tại bất kỳ thời điểm nào là tối
đa $\alpha$, nghĩa là với xác suất $1 - \alpha$, $\phi$ bằng dãy tất cả các số
không. Đây cũng được gọi là các kiểm định một phía (one-sided), kiểm định mở
(open-ended) hoặc kiểm định có công suất bằng 1 (power-one).

Không mất tính tổng quát, chúng ta có thể giả định rằng nếu $\phi_t = 1$, thì
$\phi_s = 1$ với mọi $s \geq t$. Hơn nữa, chúng ta định nghĩa $\phi_\infty =
\lim_{t \to \infty} \phi_t = \sup_{t \in \mathbb{N}} \phi_t$, nghĩa là nếu một
kiểm định không bác bỏ tại bất kỳ thời điểm hữu hạn nào, thì nó cũng không bác
bỏ tại vô cực.

Khi đó, chúng ta có thể dễ dàng xác định một \gls{st} bởi thời điểm dừng
$\tau := \inf\{t \geq 1 : \phi_t = 1\}$ với quy ước $\inf \emptyset = \infty$.
Khi đó, một \gls{test} mức ý nghĩa $\alpha$ yêu cầu $\tau < \infty$ với xác
suất tối đa $\alpha$ dưới mỗi $P \in \mathcal{P}$. Điều này có thể được diễn
đạt lại là:
\begin{equation}
    \expect_P[\phi_\tau] \leq \alpha,
\end{equation}
với mọi $P \in \mathcal{P}$ và mọi $\tau \in \mathcal{T}$.

Trong bối cảnh thời gian rời rạc, một tiến trình khả tích $M$ là một martingale
cho $P$ nếu:
\begin{equation}
    \expect_P[M_t | \mathcal{F}_{t-1}] = M_{t-1}
\end{equation}
với mọi $t \in \mathcal{T}_+$, trong đó tất cả các mối quan hệ giữa các biến
ngẫu nhiên được hiểu là đúng với xác suất 1. $M$ là một supermartingale cho $P$
nếu nó thỏa mãn điều kiện trên với dấu ``='' được nới lỏng thành ``$\leq$''.

\begin{defn}
  \begin{enumerate}[label=(\roman*)]
    \item Một tiến trình $M = (M_t)_{t \in T}$ được gọi là một test
      (super)martingale cho $\mathcal{P}$ nếu với mọi $P \in \mathcal{P}$, nó
      thỏa mãn: (a) $M$ không âm với xác suất 1 ($P$-almost surely), (b) $M$ là
      một (super)martingale dưới $P$, và (c) $\expect_P[M_0] \leq 1$. Một họ
      các tiến trình $(M^P)_{P \in \mathcal{P}}$ được gọi là một test
      (super)martingale family nếu mỗi $M^P$ là một test (super)martingale cho
      $P$.

    \item Một dãy các e-values $E = (E_t)_{t \in T}$ thích nghi với
      $\mathcal{F}$ được gọi là một e-process nếu nó thỏa mãn một trong hai
      điều kiện tương đương: (a) $\expect_P[E_\tau] \leq 1$ với bất kỳ thời
      điểm dừng $\tau \in \mathcal{T}$ và bất kỳ $P \in \mathcal{P}$; (b) tồn
      tại một họ test martingale $(M^P)_{P \in \mathcal{P}}$ sao cho $E \leq
      M^P$ với xác suất 1 ($P$-almost surely) cho mỗi $P \in \mathcal{P}$.
  \end{enumerate}
\end{defn}

\subsection{Bất đẳng thức Ville}

\begin{thm}[Bất đẳng thức Ville]
  Nếu $M$ là một e-process cho $\mathcal{P}$, thì ba phát biểu tương đương sau
  đây sẽ đúng: \begin{equation}
    P\biggl( \exists t \in \mathbb{N} : M_t \geq \frac{1}{\alpha} \biggr)
    \leq \alpha \text{ với mọi } P \in \mathcal{P} \text{ và } \alpha \in [0,
    1].
  \end{equation} trong đó $\tau$ là bất kỳ thời điểm dừng $\mathcal{F}$ nào
  nhận các giá trị trong $[0, \infty]$.
\end{thm}

\subsection{Confidence interval, Confidence sequence}

\begin{defn}
  \begin{enumerate}[label=(\roman*)]
    \item Với $\alpha \in [0, 1]$, một tập $C(\alpha)$ được gọi là $(1 -
      \alpha)$-confidence interval cho hàm $\vartheta$ nếu $\prob_P(\vartheta)
      \in C(\alpha) ) \ge 1 - \alpha$ với mọi $P \in \mathcal{P}$.
    \item Với $\alpha \in [0, 1]$, một $(1 - \alpha)$-\gls{cs} cho
      một hàm $\vartheta$ là một tiến trình các tập hợp $(C_t(\alpha))_{t \in
      \mathbb{N}}$ sao cho $C_\tau(\alpha) \subseteq \Theta$ là một $(1 -
      \alpha)$-CI (khoảng tin cậy) với bất kỳ thời điểm dừng $\tau$ nào. Nó
      cũng có định nghĩa tương đương sau đây: với mọi $P \in \mathcal{P}$,
      chúng ta có: \begin{equation}
      \prob(\vartheta(P) \in C_t(\alpha) \text{ với mọi } t \in \mathbb{N})
        \geq 1 - \alpha.
      \end{equation}
  \end{enumerate}
\end{defn}

\subsection{Asymptotic confidence sequence}

\begin{defn}
  $(\hat\mu_t \pm \bar{B}_t)_{t \ge 1}$ được gọi là $(1 - \alpha)$-\gls{acs}
  cho $\mu$ nếu tồn tại một $(1 - \alpha)$-\gls{cs} cho $\mu$ được
  cho bởi $(\hat\mu_t \pm \bar{B}_t^\ast)_{t \ge 1}$ và
  $\bar{B}_t^\ast / \bar{B}_t \stackrel{a.s}{\to} 1$.
\end{defn}

\subsection{Kiểm định giả thuyết}

Ta có mệnh đề sau:

\begin{prop}
  Mọi level-$\alpha$ sequential test $\phi$ cho $\mathcal{P}$ đều có thể nhận
  được bằng cách lấy ngưỡng $1/\alpha$ một e-process $M$ nào đó cho
  $\mathcal{P}$.
\end{prop}

Do đó, các e-process cung cấp một khung lý thuyết hoàn chỉnh cho kiểm định
tuần tự. Thay vì tập trung vào chính kiểm định đó, chúng ta có thể tập trung
vào việc xây dựng các e-process. Tuy nhiên, chúng ta không chỉ coi e-processes
là một công cụ để suy ra các kiểm định. Chúng ta coi e-process là các thước đo
bằng chứng chống lại giả thuyết không: giá trị của chúng càng lớn (hoặc tăng
lên), chúng ta càng có nhiều bằng chứng chống lại giả thuyết không. Chúng ta
có thể đặt ngưỡng cho chúng để đưa ra quyết định nếu muốn, nhưng chúng vẫn giữ
được khả năng diễn giải ngay cả khi không cần đặt ngưỡng.

\subsubsection{Kiểm định qua cá cược}

Có một cách đơn giản để xem các test (super)martingale như các tiến trình tài
sản phát sinh từ các trò chơi. Ý tưởng chủ chốt như sau:

Để kiểm định một giả thuyết không $\mathcal{P}$, chúng ta thiết lập một trò
chơi may rủi, trong đó nhà thống kê  thực hiện đặt cược vào mỗi quan sát trước
khi quan sát nó. Trò chơi này phải thỏa mãn hai tính chất:
\begin{enumerate}[label=(\alph*)]
  \item Nếu giả thuyết không là đúng, nghĩa là $P \in \mathcal{P}$, thì không
    một chiến lược cá cược nào có thể kiếm tiền một cách hệ thống (tức là, nó
    có thể kiếm tiền do may mắn, nhưng không thể đảm bảo kiếm được tiền);

  \item Nếu giả thuyết không là sai, thì phải tồn tại các chiến lược cá cược
    ``tốt'' có thể kiếm được tiền (mặc dù người ta có thể không bao giờ chắc
    chắn về lợi nhuận trong ngắn hạn, nhưng có một cách để đảm bảo kiếm được
    tiền trong dài hạn).
\end{enumerate} Nếu một trò chơi như vậy có thể được thiết kế, thì tài sản của
nhà thống kê trong trò chơi đó trực tiếp là một thước đo bằng chứng chống lại
giả thuyết không: tài sản càng nhiều (hoặc chính xác hơn là tỉ số giữa tài sản
cuối cùng và tài sản ban đầu càng lớn), thì càng có nhiều bằng chứng cho thấy
giả thuyết không có khả năng là sai.

Chúng ta hãy bắt đầu với ví dụ đơn giản nhất, như đã thảo luận trong phần
trước. Chúng ta chỉ ra cách diễn giải việc kiểm định $P$ trong
\algorithmname~\ref{alg:simple-P}.

\begin{algorithm}[htbp]
  \caption{Kiểm định qua cá cược: trường hợp $P$ đơn}
  \label{alg:simple-P}
  Khởi tạo tài sản ban đầu cho nhà thống kê $W_0 = 1$\;
  \For{$t = 1$, $2$, $\ldots\,$}{
    Nhà thống kê tuyên bố một khoản cược $S_t \colon \mathcal{X} \to [0,
    \infty]$ sao cho \begin{equation*}
      \expect_P[S_t(X) \mid X_1, \dots, X_{t - 1}] \le 1
    \end{equation*}\vspace{-1\baselineskip}\;
    Quan sát $X_t$\;
    Cập nhật tài sản của nhà thống kê $W_t = W_{t - 1} \cdot S_t(X_t)$\;
  }
\end{algorithm}

Có ba điểm đáng chú ý cần bình luận ngay. Thứ nhất, trong trò chơi trên, ràng
buộc đối với $S_t$ đơn giản là nó là một e-variable cho $\mathcal{P}$ có điều
kiện dựa trên quá khứ. Điều này hàm ý ngay lập tức rằng đối với bất kỳ chiến
lược cá cược nào, tiến trình tài sản $W$ là một test supermartingale cho
$\mathcal{P}$.

Thứ hai, người ta có thể hỏi: đâu là khoản cược tối ưu? Giả sử để đơn giản rằng
giả thuyết thay thế $Q$ là đơn giản, và $P \ll Q$. Khi đó, tài sản có thể tăng
trưởng theo hàm mũ dưới $Q$, và do đó việc tối ưu hóa số mũ là tự nhiên. Điều
này có nghĩa là chọn $S_t$ để tối ưu hóa $\expect_Q[\log S_t(X) | X_1, \dots,
X_{t-1}]$ với điều kiện $S_t$ là một e-variable cho $\mathcal{P}$. Câu trả lời
đã được rút ra, đó là tỉ số khả năng (likelihood ratio) của $Q$ so với $P$, có
điều kiện dựa trên $X_1, \dots, X_{t-1}$ (phép điều kiện này không liên quan
nếu $P$ và $Q$ là các phân phối tích i.i.d trên $\mathcal{X}^\infty$, nhưng nó
lại quan trọng nếu chúng không phải như vậy).

Điểm cuối cùng là nếu người ta muốn kiểm định một giả thuyết không hợp
$\mathcal{P}$, thay đổi duy nhất đối với giao thức là yêu cầu $S_t$ phải là một
e-variable cho $\mathcal{P}$, nghĩa là ràng buộc trong giao thức trên phải
giữ nguyên với mọi $P \in \mathcal{P}$. Khi đó, tiến trình tài sản $W$ là một
testsupermartingale cho $\mathcal{P}$.

Bây giờ, chúng ta có thể tự nhiên thắc mắc làm thế nào các e-process có thể
được diễn giải trong khuôn khổ kiểm định qua cá cược này, trái ngược với việc
chỉ đơn thuần là các test supermartingale. Câu trả lời là diễn giải lý thuyết
trò chơi của chúng phức tạp hơn một chút. Nhà thống kê không chơi một trò chơi
duy nhất chống lại $\mathcal{P}$, mà thay vào đó chơi một trò chơi riêng biệt
chống lại mỗi $P \in \mathcal{P}$, và chấp nhận tài sản ròng của mình là tài
sản tệ nhất trong tất cả các trò chơi này. Điều này được làm rõ trong
\algorithmname\ref{alg:composite-P}.

\begin{algorithm}[htbp]
  \caption{Kiểm định qua cá cược: trường hợp $P$ hợp}
  \label{alg:composite-P}

  \For{$P \in \mathcal{P}$}{
    Khởi tạo tài sản ban đầu của nhà thống kê cho trò chơi $P$ $W_0^P = 1$\;
  }
  \BlankLine
  \For{$t = 1$, $2$, $\ldots\,$}{
    Nhà thống kê tuyên bố một khoản cược $S^P_t \colon \mathcal{X} \to [0,
    \infty]$ sao cho \begin{equation*}
      \expect_P[S^P_t(X) \mid X_1, \dots, X_{t - 1}] \le 1
    \end{equation*}\vspace{-1\baselineskip}\;
    Quan sát $X_t$\;
    Cập nhật tài sản trong $P$ của nhà thống kê $W^P_t = W^P_{t - 1} \cdot
    S^P_t(X_t)$\;

    Tài sản tổng thể của nhà thống kê được cập nhật $W_t = \inf_{P \in
    \mathcal{P}} W^P_t$\;
  }
\end{algorithm}

\subsubsection{Trường hợp giả thuyết không và thay thế đều đơn}

Giả sử chúng ta đang kiểm định một giả thuyết đơn $P$ đối với một giả thuyết
đơn $Q \ll P$. Nếu chúng ta quan sát dữ liệu i.i.d $X_1, X_2, \dots$ một cách
tuần tự, thì tiến trình tỉ số khả năng $M$ được cho bởi:
\begin{equation}
    M_0 = 1 \text{ và } M_t = \prod_{i=1}^{t} \frac{dQ}{dP}(X_i) \text{ với } t
    \in \mathbb{N},
\end{equation}
là một e-process thích nghi với sự lọc được tạo ra bởi dữ liệu, trong đó
$dQ/dP(X)$ là tỉ số khả năng của $Q$ đối với $P$ khi quan sát $X$. Chúng ta đã
bắt gặp e-process này trong Mục 7.2, được sử dụng trong SPRT. Hơn nữa, dễ dàng
thấy rằng $M$ là một test martingale cho $P$. Trên thực tế, ngay cả khi dữ liệu
không phải i.i.d, tỉ số khả năng dựa trên $n$ điểm dữ liệu đầu tiên:
\begin{equation}
    M_n = \frac{dQ}{dP}(X_1, \dots, X_n),
\end{equation}
vẫn là một test martingale cho $P$ (có thể kiểm tra bằng phép lấy điều kiện).

Khi kiểm định một giả thuyết không đơn giản $P$, việc sử dụng một test
martingale để xây dựng các e-process là tối ưu: Bất kỳ e-process nào cho $P$
đều có thể bị chặn trên bởi Snell envelope của nó dưới $P$ (vốn là một test
supermartingale), và đến lượt nó, lại có thể bị chi phối bởi test martingale
cho $P$ được cho bởi phân tích Doob.

Trong bối cảnh kiểm định giả thuyết không đơn $P$ đối với $Q \ll P$, tỉ số
khả năng $M_t$ là tối ưu logarit trong số tất cả các e-variable sử dụng $t$
điểm dữ liệu đầu tiên với mọi $t \in \mathbb{N}$.

Khi đó, khoản cược log-optimal là \begin{equation}
  S_t(X) = \frac{dQ}{dP}(X) = \frac{q(X)}{p(X)},
\end{equation} trong đó có đẳng thức sau nếu $P$, $Q$ các các hàm mật độ xác
suất $p$ và $q$ tương ứng. Khi đó $\expect_P[W_\tau] \le 1$, $\expect_P[\log
W_\tau] \le 0$. Trong khi đó, $\expect_Q[\log W_T] > 0$ được lấy tối đa và
tương đương $\KL(Q, P)$.

\subsubsection{Các trường hợp khác}

Các trường hợp giả thuyết không hợp, giả thuyết thay thế đơn; giả thuyết
không đơn, giả thuyết thay thế hợp; hay cả giả thuyết không và thay thế đều hợp
đều khá phức tạp. Do đó, báo cáo này không trình bày các phần này. Ta có thể
tham khảo thêm trong \cite{ramdas2023savi}
hoặc \cite{ramdas2025ebook}.


\subsection{Ước lượng}

\subsubsection{Khoảng tin cậy cho giá trị trung bình}

\begin{exmp}
  Nếu $X_1$, $X_2$, $\ldots\,$ là các biến ngẫu nhiên i.i.d Gaussian hoặc bị
  chặn trong $[-1, 1]$ thì \begin{equation}
    \frac{ \sum_{i = 1}^n X_i }{ n } \pm 1.71 \sqrt{ \frac{ \log\log(2n) + 0.72
    \log ( 10.4 / \alpha ) }{ n } }
  \end{equation} là một $(1 - \alpha)$ confidence sequence cho $\mu$, tương tự
  như \begin{equation}
    \bigcap_{s \le n} \frac{ \sum_{i = 1}^n X_i }{ n } \pm 1.71 \sqrt{ \frac{
      \log\log(2n) + 0.72 \log ( 10.4 / \alpha ) }{ n } }.
  \end{equation}
  Ta thấy rằng, độ rộng của $(1 - \alpha)$ confidence sequence lớn hơn so với
  khoảng tin cậy tương ứng nhưng chênh lệch không thực sự ảnh hưởng nhiều. Để
  thấy sự khác biệt giữa hai kết quả này. Ta dùng đến tỉ lệ bao phủ sai tích
  lũy (cummulative miscoverage rate). Tỉ lệ này ứng với mỗi $t$ chính là tỉ lệ
  tồn $t$ khoảng tin cậy đầu tiên bị bao phủ sai với các lần chạy. Do khoảng
  tin cậy chỉ đảm bảo sai lầm ở các lần chạy ở một thời điểm $t$ là dưới
  $\alpha$ nên khi xét đồng thời $t$ khoảng tin cậy thì tỉ lệ bao phủ đồng thời
  của $t$ khoảng tin cậy gần như bằng $0$ (hay tỉ lệ bao phủ sai gần như là
  $1$). Tuy nhiên, theo định nghĩa của confidence sequence, tỉ lệ bao phủ sai
  tích lũy luôn được đảm bảo dưới $\alpha$ do tính đồng thời của nó. Điều này
  cũng cho thấy tính hợp lệ tại mọi thời điểm của confidence sequence so với
  khoảng tin cậy. (Xem minh họa tại trang số 8, nửa bài giảng sau.)
\end{exmp}

\subsubsection{Khoảng tin cậy cho phân vị}

\paragraph{Một phân vị cố định}

\begin{exmp}
  Đặt \begin{equation*}
    u_t := \sqrt{ \frac{ 0.73 \log\log(2.04t) + 0.52 \log(9.97/\alpha) } { t }
    }
  \end{equation*} Khi đó, \begin{equation*}
    \prob(\forall t \in \mathbb{N} : \hat{Q}_t(1/ 2 - u_t) \le Q(1/2) \le
    \hat{Q}_t(1/2 + u_t)) \ge 1 - \alpha.
  \end{equation*} Tức là \begin{equation*}
    [\hat{Q}_t(1/ 2 - u_t), \hat{Q}_t(1/2 + u_t))], \quad \forall t
  \end{equation*} là một $(1 - \alpha)$ confidence sequence cho phân vị $1/2$
(trung vị).

\end{exmp}


\paragraph{Toàn bộ phân vị đồng thời}

\begin{exmp}
  Đặt \begin{equation*}
    u_t := \sqrt{ \frac{ \log\log(et) + 0.75 \log(34 / \alpha) } { t } }.
  \end{equation*} Khi đó, \begin{equation*}
    \prob(\forall t \in \mathbb{N}, p \in (0, 1): \hat{Q}_t(p - u_t) \le Q(p)
    \le \hat{Q}_t(p + u_t) ) \ge 1 - \alpha.
  \end{equation*} Tức là ta có đồng thời các khoảng tin cậy cho từng thời điểm
  với cùng độ rộng cho tất cả phân vị khác nhau. Và chúng đều tạo thành một $(1
  - \alpha)$ confidence sequence đồng thời cho tất cả các phân vị.
\end{exmp}

\subsubsection{Ước lượng ma trận hiệp phương sai tuần tự}

Xét biến ngẫu nhiên $X \in \mathbb{R}^d$ với kỳ vọng $\expect[X] = 0$ và điều
kiện bị chặn $|X_i| \le b$. Sai số hội tụ của ước lượng $\hat{\Sigma}_n$ so với
ma trận thực $\Sigma$ được chặn trên với xác suất cao như sau:
\begin{equation}
  \big\| \hat{\Sigma}_n - \Sigma \big\|_{\text{op}} \lesssim \sqrt{\frac{b^2
  \log(d \log n)}{n}} + \frac{b^2 \log(d \log n)}{n}.
\end{equation}

\subsubsection{Ước lượng qua cá cược}

Cho $X_1$, $X_2$, $\ldots\,$ là các biến ngẫu nhiên i.i.d thuộc $[0, 1]$ với
trung bình $\mu$. Ta thiết kế trò chơi như
\algorithmname~\ref{alg:esitmate-game}. Khi đó đặt \begin{equation*}
  C_t := \{ m \in [0, 1]: K^m_t < 1/\alpha \}
\end{equation*} với ý nghĩa rằng nhà khoa học không thắng được nhiều tiền. Ta
có định lý (có thể chứng minh thông qua kiểm định qua cá cược) nói rằng
$(C_t)_{t \ge 1}$ là một confidence sequence cho $\mu$ với mọi chiến lược đặt
cược.

\begin{algorithm}[htbp]
  \caption{Ước lượng qua cá cược}
  \label{alg:esitmate-game}

  Khởi tạo vốn ban đầu $K^m_0 = 1$ cho mỗi trò chơi $m \in [0, 1]$\;

  \For{$t = 1, 2, \dots$}{
    Với mỗi $m \in [0, 1]$, nhà thống kê đặt cược $\lambda_t^m \in [ -1/( 1- m
    ), 1/m ]$\;
    Quan sát $X_t$\;
    Cập nhật tài sản trong game $m$ thành $K^m_t = K^m_{t - 1} \cdot (1 +
    \lambda^m_t(X_t - m))$\;
  }
\end{algorithm}

Khi đó, vấn đề trở thành ta cần một chiến lược đặt cược để các tập hợp này càng
nhỏ càng tốt. Một trong những chiến lược là GRAPA (Grow Rate Adaptive to the
Particular Alternative), ta đặt cược như sau: \begin{equation}
  \lambda^m_t(P) := \arg\max_{\lambda \in [-1/(1 - m), 1/m]} \expect_P[\log(1 +
  \lambda(X_t - m)) \mid \mathcal{F}_{t - 1}].
\end{equation} Nhưng ta không biết $P$, lời giải xấp xỉ: lấy đạo hàm theo biến
$\lambda$, cho đạo hàm bằng 0 (áp dụng điều kiện KKT), khai triển Taylor và
thay thế bằng các ước lượng thực nghiệm. Khi đó ta có giá trị cho khoản cược là
\begin{equation*}
  \lambda^m_t = \frac{\hat\mu_t - m}{ \hat\sigma_t^2 + (\hat\mu_t - m)^2 },
\end{equation*} trong đó các $\mu_t$ và $\sigma_t^2$ được tính từ $t - 1$ mẫu
đầu.

\subsection*{Chiến lược cá cược 2: Phương pháp hỗn hợp (Mixture)}

Chiến lược thứ hai là Mixture. Chiến lược này được củng cố và phát triển thông
qua hàng loạt các nghiên cứu gần đây, tiêu biểu như Waudby-Smith \& Ramdas
(2024), Jun \& Orabona (2024), và Shekhar \& Ramdas (2024). Theo đó, confidence
sequence được định nghĩa dựa trên một tích phân hỗn hợp như sau:
\begin{equation}
  C_t := \Biggl\{ m \in [0,1] : \int_{-1/(1-m)}^{1/m} \prod_{i=1}^t \big(1 +
  \lambda(X_i - m)\big) \, d\nu(\lambda) < \frac{1}{\alpha} \Biggr\}
\end{equation} Về mặt triển khai thực tiễn, việc tính toán tích phân liên tục
thường được thay thế bằng phương pháp rời rạc hóa (discretize) không gian của
hỗn hợp. Khía cạnh đáng chú ý là phép xấp xỉ này hoàn toàn không làm ảnh hưởng
đến tính hợp lệ thống kê (validity) của confidence sequence.

Tuy nhiên, nhằm bảo toàn sức mạnh kiểm định (power) khi lượng dữ liệu tăng lên,
lưới rời rạc hóa của hỗn hợp bắt buộc phải được tinh chỉnh mịn hơn (finer) theo
thời gian. Cuối cùng, để đánh giá và phân tích hiệu năng của chiến lược này, ta
có thể thiết lập mối liên hệ trực tiếp với bài toán Tối ưu hóa Danh mục Đầu tư
của Cover (Cover's Portfolio Optimization). Dựa trên cơ sở lý thuyết đó, việc
sử dụng một phân phối hỗn hợp đều (uniform mixture) thường là phương án đủ để
đảm bảo tính tối ưu.

\subsubsection{}

Trong suy luận nhân quả và phân tích thực nghiệm, sự kết hợp giữa Tác động can
thiệp trung bình (ATE) và asymptotic confidence sequence (AsympCS) mang lại một
giải pháp đột phá cho bài toán đánh giá dữ liệu liên tục.

Trong khi ATE đóng vai trò là thước đo cốt lõi giúp xác định hiệu quả trung
bình thực sự của một sự can thiệp (như một tính năng công nghệ mới hay một phác
đồ điều trị) trên toàn bộ quần thể, thì việc liên tục theo dõi ATE để ra quyết
định sớm lại rất dễ vi phạm các nguyên tắc thống kê truyền thống. Đây là lúc
AsympCS trở thành mảnh ghép hoàn hảo. Bằng cách nới lỏng các ràng buộc toán học
khắt khe tuyệt đối thành dạng xấp xỉ tiệm cận (thường dựa trên Định lý giới hạn
trung tâm), AsympCS cung cấp một dải an toàn cho phép bạn liên tục "nhìn" vào
kết quả ATE theo thời gian thực (anytime-valid inference). Nhờ đó, các nhà
nghiên cứu và kỹ sư dữ liệu có thể linh hoạt dừng thử nghiệm sớm ngay khi hiệu
quả của can thiệp đã đủ rõ ràng, vừa tiết kiệm thời gian, chi phí, lại vừa đảm
bảo tính hợp lệ thống kê chặt chẽ trong dài hạn.

